\documentclass[10pt,a4paper]{amsart}
\usepackage[margin=19mm]{geometry}
\usepackage{amsmath,amssymb,mathtools}
\usepackage[scr=rsfs,cal=euler]{mathalfa}
\usepackage{xfrac}
\usepackage[unicode,hidelinks,bookmarks=false]{hyperref}
\newcommand{\ie}{i.e.\ }
\newcommand{\cf}{cf.\ }
\newcommand{\resp}{respectively\ }
\newcommand{\Subsection}{\subsection}
\providecommand{\nobreakdash}{\nobreak\hbox{-}}
\setlength{\emergencystretch}{2em}
\multlinegap0pt
\usepackage{bm}
\usepackage{bbm}
\usepackage{dsfont}
\usepackage{xspace}
\usepackage{cancel}
\usepackage{mathtools}
\usepackage{amsthm}
\makeatletter
\@ifundefined{theorem}{\newtheorem{theorem}{Theorem}}{}
\@ifundefined{lemma}{\newtheorem{lemma}[theorem]{Lemma}}{}
\makeatother
\DeclareMathOperator{\nil}{Nil}
\DeclareMathOperator{\spa}{span}
\title[A Frattini theory for evolution algebras]{A Frattini theory for evolution algebras}
\author{Manuel Ladra, Andr\'es P\'erez-Rodr\'iguez}
\date{Source: arXiv:2507.01935v1 (CC BY 4.0)}
\begin{document}
\maketitle

\noindent\emph{Source and licence.} A Frattini theory for evolution algebras. Version arXiv:2507.01935v1 (CC BY 4.0), distributed under the Creative Commons Attribution 4.0 International licence, as recorded in the arXiv licence metadata for this exact version and shown on the abstract page https://arxiv.org/abs/2507.01935v1. Full rights notice: licenses/arxiv-2507.01935-CC-BY-4.0.txt.\par\smallskip

\noindent\textit{Editorial scope.} Counted unit: the Theorem whose source label is \texttt{th:nilrad\_dim1} in the article (source0.tex lines 288--293, source label \texttt{th:nilrad\_dim1}), delivered together with the article's own complete original proof of it (source0.tex lines 294--318, one \texttt{proof} environment of 25 lines). No mathematics is written, repaired or re-derived: statement and proof are the article's own text line for line. Required context retained uncounted: lem:lem1 (lines 280--282). Every definition, display and label of those lines is retained unchanged. Omitted from this package and not redistributed: 1--279, 283--287, 319--982. Nothing inside a retained excerpt is omitted or altered: every retained body is delivered exactly as the article's lines are written, so no omission receipt of any kind accompanies this package. Citations: the retained excerpts carry no citation command at all, so no bibliography entry of the article is transcribed and no reference is redistributed. Reference adaptation: none was needed and none was made; every reference inside the delivered unit resolves inside it, so no unresolved reference or citation is produced. Printed slips kept literal: no printed word, symbol, hypothesis or number of the counted statement or of its proof was altered. Layout and class: the article's own class is \texttt{amsart} with packages including amssymb, a4wide, amsmath, mathtools, xy, latexsym, pdfsync, xcolor, graphicx, multirow, array, upgreek, pstricks-add, enumerate; the wrapper is the lane's pinned \texttt{amsart} editorial wrapper at 10pt on A4, which restates the theorem-like environments the retained text needs and adds the article's own macro definitions that the retained text needs (\texttt{\textbackslash nil}, \texttt{\textbackslash spa}), with extra wrapper packages (bm, bbm, dsfont, xspace, cancel, mathtools). The delivered numbering is the unit's own. Assets: no retained span contains a figure environment, a graphics-inclusion command, a PDF-page inclusion, a table or a TikZ picture; no image file is copied into this package, the package declares no asset dependency and no image file is redistributed with it. Licence: version arXiv:2507.01935v1 of this article is distributed under the Creative Commons Attribution 4.0 International licence, as recorded in the arXiv licence metadata for this exact version and shown on the abstract page https://arxiv.org/abs/2507.01935v1. A full rights notice is delivered at licenses/arxiv-2507.01935-CC-BY-4.0.txt. Characters: every retained excerpt is pure ASCII, so the delivered engine, XeLaTeX with its default Latin Modern text font, reports no missing character.
\begin{lemma}\label{lem:lem1}
	Let $\mathcal{E}\in\mathcal{T}_\mathbb{K}$. Then, its derived subalgebra, $\mathcal{E}^2$, is contained in every maximal nilpotent ideal.
\end{lemma}

\begin{theorem}\label{th:nilrad_dim1}
	Let $\mathcal{E}=\mathcal{E}_k(\lambda_1,\dots,\lambda_n)\in\mathcal{T}_\mathbb{K}$. Then, its nilradical, $\nil(\mathcal{E})$, exists. In fact, assuming, without loss of generality, that $k\geq2$ and $\lambda_1,\dots,\lambda_k\neq0$, it holds that
	\begin{align*}
	\nil\big(\mathcal{E}_k(\lambda_1,\dots,\lambda_n)\big)=\spa\{\lambda_2e_1-\lambda_1e_2,\lambda_3e_1-\lambda_1e_3,\dots,\lambda_ke_1-\lambda_1e_k,e_{k+1},\dots,e_n\}.
	\end{align*}
\end{theorem}

\begin{proof}
	For simplicity, we denote by $\mathcal{N}$ the subspace described just above and prove that it is the only maximal nilpotent ideal of  $\mathcal{E}_k(\lambda_1,\dots,\lambda_n)$ with $k\geq2$ and $\lambda_1,\dots,\lambda_k\neq0$. First, $\mathcal{N}$ is an ideal because $\mathcal{E}^2=\spa\{e_1+\dots+e_k\}\subset\mathcal{N}$. In fact, we have that
	\begin{align*}
		(\lambda_2e_1-\lambda_1e_2)&+(\lambda_3e_1-\lambda_1e_2)+\dots+(\lambda_ke_1-\lambda_1e_k)\\&
		=(\lambda_2+\lambda_3+\dots+\lambda_k)e_1-\lambda _1e_2-\lambda_1e_3-\dots-\lambda_1e_k\\&
		=-\lambda_1e_1-\lambda _1e_2-\lambda_1e_3-\dots-\lambda_1e_k
		=-\lambda_1(e_1+\dots+e_k);
	\end{align*}
	and, as $\lambda_1\neq0$ by hypothesis, we conclude that $e_1+\dots+e_k\in\mathcal{N}$. Secondly, $\mathcal{N}$ is maximal since it has codimension one. Thirdly, $\mathcal{N}$ is nilpotent since $\mathcal{N}^3=\mathcal{N}^2\mathcal{N}=0$. In fact, it holds that
	\begin{align}\label{eq:para_siguiente_th}
		(\lambda_ie_1-\lambda_1e_i)(e_1+\dots+e_k)=\lambda_i\lambda_1(e_1+\dots+e_k)-\lambda_1\lambda_i(e_1+\dots+e_k)=0,
	\end{align}
for any $2\leq i\leq k$ and $e_i(e_1+\dots+e_k)=0$ for any $k+1\leq i\leq n$. 

Next, for the sake of contradiction, suppose that there exists another maximal nilpotent ideal $\mathcal{M}$ and consider a nonzero element $u=\sum_{i=1}^n\mu_ie_i$ such that $u\in\mathcal{M}$ but $u\notin\mathcal{N}$. Then, we have that 
\begin{multline}\label{op}
	u+\frac{\mu_2}{\lambda_1}(\lambda_2e_1-\lambda_1e_2)+\dots+\frac{\mu_k}{\lambda_1}(\lambda_ke_1-\lambda_1e_k)-\mu_{k+1}e_{k+1}-\dots-\mu_ne_n\\=\left(\mu_1+\frac{\mu_2\lambda_2}{\lambda_1}+\dots+\frac{\mu_k\lambda_k}{\lambda_1}\right)e_1\neq0.
\end{multline}
Note that if \eqref{op} were equal to zero, there would be a contradiction with the fact that $u\notin\mathcal{N}$. Consequently, $\mu_1+\frac{\mu_2\lambda_2}{\lambda_1}+\dots+\frac{\mu_k\lambda_k}{\lambda_1}\neq0$. Moreover, by Lemma \ref{lem:lem1}, it holds that $e_1+\dots+e_k\in\mathcal{M}$. Additionally, it holds that $u(e_1+\dots+e_k)=0$. Otherwise, $u(e_1+\dots+e_k)=\mathbb{K}^*(e_1+\dots+e_k)$ and, consequently, $\mathcal{M}^{\langle n\rangle}\neq0$ for any $n\in\mathbb{N}$, which contradicts the nilpotency of $\mathcal{M}$. Finally, putting all this together, we get that 
\begin{align*}
	0=&\left(u+\frac{\mu_2}{\lambda_1}(\lambda_2e_1-\lambda_1e_2)+\dots+\frac{\mu_k}{\lambda_1}(\lambda_ke_1-\lambda_1e_k)-\mu_{k+1}e_{k+1}-\dots-\mu_ne_n\right)(e_1+\dots+e_k)\\=&\left(\mu_1+\frac{\mu_2\lambda_2}{\lambda_1}+\dots+\frac{\mu_k\lambda_k}{\lambda_1}\right)e_1(e_1+\dots+e_k)
	=\left(\mu_1+\frac{\mu_2\lambda_2}{\lambda_1}+\dots+\frac{\mu_k\lambda_k}{\lambda_1}\right)e_1^2,
\end{align*}
which is impossible, thus leading to a contradiction with the assumption that $\mathcal{M}$ is another maximal nilpotent ideal.
\end{proof}

\end{document}
